\section{Applications to Markov chains}
We use column probability vectors. Entry $p_{ij}$ is the probability of
moving \emph{from state $j$ to state $i$}, so each column of the transition
matrix sums to $1$ and its entries are nonnegative. The distribution evolves
by $x_{k+1}=Px_k$. A steady-state vector is a probability vector $q$ with
$Pq=q$; solve $(P-I)q=0$ together with $\sum_iq_i=1$.
A stochastic matrix is regular if some positive integer power has every entry
strictly positive. For a regular matrix, Theorem 18 guarantees a unique
steady-state vector and convergence $P^kx_0\to q$ for every initial probability vector.

\practice{4.9.P1}{A city resident's next two locations}
The migration matrix from Example 1 is
$M=\left(\begin{smallmatrix}.95&.03\\.05&.97\end{smallmatrix}\right)$,
with city first and suburbs second. Starting from $x_0=(1,0)$,
$x_1=Mx_0=(.95,.05)$. Thus (a) the probability of living in the suburbs
next year is $5\%$. A second multiplication gives
$x_2=(.95^2+.03(.05),.05(.95)+.97(.05))=(.904,.096)$.
Thus (b) the probability after two years is $9.6\%$.

\practice{4.9.P2}{Test the steady-state equation rather than its appearance}
For the given $P,q$, direct multiplication gives
\[
 \begin{pmatrix}.6&.2\\.4&.8\end{pmatrix}
 \begin{pmatrix}.3\\.7\end{pmatrix}
 =\begin{pmatrix}.18+.14\\.12+.56\end{pmatrix}
 =\begin{pmatrix}.32\\.68\end{pmatrix}\ne q.
\]
Although $q$ is a probability vector, it is not a steady-state vector.

\practice{4.9.P3}{A long-run migration percentage}
Solve $Mq=q$: the first equation reduces to $.05q_1=.03q_2$.
Together with $q_1+q_2=1$, this gives $q=(3/8,5/8)$.
Every entry of $M$ is positive, so it is regular; Theorem 18 makes this
the limiting distribution. Thus $62.5\%$ will live in the suburbs after many years.

\exercise{4.9.1}{Radio listeners at two station breaks}
\step{(a) Build columns from current stations.}
Order the states as news, music for both rows (next station) and columns
(current station). The probabilities give
\[
 P=\begin{pmatrix}.7&.6\\.3&.4\end{pmatrix}.
\]
Column one describes a news listener; column two describes a music listener.
Each column sums to $1$.
\step{(b) Use the stated initial audience.}
Everyone initially listens to news, so $x_0=(1,0)$.
\step{(c) Count the breaks, not the elapsed minutes.}
At 9:25 there have been two breaks, at 8:30 and 9:00.
Then $x_1=(.7,.3)$ and
$x_2=(.7(.7)+.6(.3),.3(.7)+.4(.3))=(.67,.33)$.
Thus $33\%$ will be listening to music.

\exercise{4.9.2}{Three equally treated food choices}
\step{(a) Use foods 1, 2, 3 in both state lists.}
The probability of retaining the current food is $.6$; the other two
possibilities each have probability $.2$. Thus
\[
 P=\begin{pmatrix}.6&.2&.2\\.2&.6&.2\\.2&.2&.6\end{pmatrix}.
\]
\step{(b) Apply two transitions after the initial trial.}
Starting at food 1 gives $x_0=(1,0,0)$ and $x_1=(.6,.2,.2)$.
The second component of $x_2$ is
$.2(.6)+.6(.2)+.2(.2)=.12+.12+.04=.28$.
The requested probability is $28\%$; a single multiplication would answer
the first trial after the initial trial, not the second.

\exercise{4.9.3}{Healthy and ill students}
\step{(a) Define the state order and fill the columns.}
Use healthy, ill. A healthy student becomes ill with probability $.05$;
an ill student becomes healthy with probability $1-.55=.45$. Therefore
\[
 P=\begin{pmatrix}.95&.45\\.05&.55\end{pmatrix}.
\]
\step{(b) Advance Monday's probability vector twice.}
Monday's $x_0=(.8,.2)$ gives Tuesday
$x_1=(.95(.8)+.45(.2),.05(.8)+.55(.2))=(.85,.15)$.
Wednesday gives $x_2=(.875,.125)$. The ill percentages are $15\%$ and $12.5\%$.
\step{(c) Use a different initial condition for a currently healthy student.}
From $(1,0)$, after one day the vector is $(.95,.05)$, and after two days
the healthy probability is $.95^2+.45(.05)=.925$.
Thus the requested probability is $92.5\%$.

\exercise{4.9.4}{Good, indifferent, and bad weather}
\step{(a) Supply probabilities missing from the wording.}
Use good, indifferent, bad. Each column sums to one, so the remaining
indifferent-to-bad probability is $.3$ and bad-to-bad is $.3$.
The matrix is
\[
 P=\begin{pmatrix}.4&.5&.3\\.3&.2&.4\\.3&.3&.3\end{pmatrix}.
\]
\step{(b) Tomorrow's bad-weather probability.}
Starting from $(.5,.5,0)$, its bad component after one transition is
$.3(.5)+.3(.5)+.3(0)=.3$, or $30\%$.
\step{(c) Wednesday is two transitions from Monday.}
From Monday $(0,.6,.4)$, Tuesday is $(.42,.28,.30)$.
Wednesday's good component is $.4(.42)+.5(.28)+.3(.30)=.398$,
so the answer is $39.8\%$.

\exercise{4.9.5}{Normalize a two-state balance equation}
For $P=\left(\begin{smallmatrix}.1&.5\\.9&.5\end{smallmatrix}\right)$,
$Pq=q$ reduces to $.9q_1=.5q_2$, or $q_2=(9/5)q_1$.
Impose $q_1+q_2=1$ to get $(14/5)q_1=1$.
\answer{$q=(5/14,9/14)$. Both entries are nonnegative and sum to one.}

\exercise{4.9.6}{A second two-state equilibrium}
Here the balance equation is $.6q_1=.8q_2$, so $3q_1=4q_2$.
Together with normalization this gives
\[
 q=(4/7,3/7).
\]
Substitution gives $Pq=q$, confirming the calculation.

\exercise{4.9.7}{Three-state balance with two useful equations}
From the first row, $.3q_1=.1q_2+.1q_3=.1(1-q_1)$,
so $.4q_1=.1$ and $q_1=1/4$.
From the second, $.2q_2=.2q_1+.2q_3=.2(1-q_2)$,
so $q_2=1/2$. Normalization then gives $q_3=1/4$.
\answer{$q=(1/4,1/2,1/4)$. The remaining row equation is satisfied as well.}

\exercise{4.9.8}{A steady state despite zero transition entries}
The second and third row equations give $.5q_2=.1q_3$ and
$.9q_3=.6q_1$. Thus $q_2=q_3/5$ and $q_3=2q_1/3$.
Normalizing gives
\[
 q_1(1+2/15+2/3)=1,
 \qquad q=(5/9,2/27,10/27).
\]
Substituting these fractions into the first row gives $5/9$; their sum is
one, so they form the required probability vector.

\exercise{4.9.9}{One zero can disappear under squaring}
The given matrix is stochastic, since it has nonnegative entries and unit
column sums. Its square is
\[
 P^2=\begin{pmatrix}.2&1\\.8&0\end{pmatrix}^2
 =\begin{pmatrix}.84&.2\\.16&.8\end{pmatrix}.
\]
Every entry of this power is strictly positive. Therefore $P$ is regular,
even though $P$ itself has a zero entry.

\exercise{4.9.10}{An absorbing state prevents regularity}
For $P=\left(\begin{smallmatrix}1&.3\\0&.7\end{smallmatrix}\right)$,
the first column is $(1,0)$, and $P(1,0)=(1,0)$.
Inductively the first column of every positive power remains $(1,0)$.
Thus each power has a zero entry, so $P$ is not regular.

\exercise{4.9.11}{Long-run radio preferences}
For the matrix in Exercise 4.9.1, the balance equation is $.3q_N=.6q_M$.
Thus $q_N=2q_M$, and normalization gives (a) $q=(2/3,1/3)$.
All entries of $P$ are positive, so it is regular. Theorem 18 implies
(b) that the long-run fraction listening to the news is $2/3$, regardless
of the initial audience. The finite two-break answer need not equal this limit.

\exercise{4.9.12}{No long-run preference between symmetric foods}
The matrix of Exercise 4.9.2 sends $(1/3,1/3,1/3)$ to itself; every row
also sums to one. All entries are positive, so Theorem 18 makes this the
unique limiting probability vector. Hence none of the three foods is preferred
in the long run: each is chosen with probability $1/3$.

\exercise{4.9.13}{The long-run illness rate}
For the healthy/ill matrix, $.05q_H=.45q_I$ implies $q_H=9q_I$.
Normalization gives (a) $q=(.9,.1)$.
Since this matrix has all positive entries, it is regular; thus (b) the
long-run probability that a specific student is ill is $.1$, or $10\%$.
Theorem 18 applies to every initial probability vector, including certainty
of being healthy or ill today, so today's state does not change that limit.

\exercise{4.9.14}{Long-run good weather}
The third row of the weather matrix gives $q_B=.3$ immediately because
$q_G+q_I+q_B=1$. Then $q_G+q_I=.7$ and
$q_I=.3q_G+.2q_I+.4(.3)$.
Substitute $q_I=.7-q_G$ to obtain
$.8(.7-q_G)=.3q_G+.12$, whence $q_G=.4$ and $q_I=.3$.
The matrix is strictly positive, so the long-run good-weather probability
is $40\%$.

\exercise{4.9.15}{A migration percentage in the long run}
With states California and rest of the United States, use
\[
 P=\begin{pmatrix}.9821&.0029\\.0179&.9971\end{pmatrix}.
\]
The steady-state equation gives $.0179q_{CA}=.0029q_{rest}$.
Since $q_{rest}=1-q_{CA}$,
\[
 q_{CA}=\frac{.0029}{.0179+.0029}=\frac{29}{208}\approx .13942308.
\]
All entries are positive, hence the matrix is regular. The eventual share
is approximately $13.9423\%$. The initial $11.7\%$ affects the transient
distributions, but not this limit, under the stipulated constant migration model.

\exercise{4.9.16}{Allocate a car fleet across three locations}
Order states as City Airport, Downtown, Metro Airport, using the printed
\emph{columns} as rental locations and rows as return locations:
\[
 P=\begin{pmatrix}.90&.01&.09\\.01&.90&.01\\.09&.09&.90\end{pmatrix}.
\]
The downtown balance gives $.10q_D=.01q_C+.01q_M=.01(1-q_D)$,
so $.11q_D=.01$ and $q_D=1/11$.
The full steady state is $(91/209,1/11,9/19)$; its entries sum to one and
satisfy $Pq=q$. Because the matrix is positive, it is regular and this
steady distribution describes a typical long-run day under the model.
For about $2000$ cars, downtown gets $2000/11\approx181.82$ cars,
or approximately \textbf{182 cars} rented or ready to rent there.

\exercise{4.9.17}{A stochastic matrix has eigenvalue one}
\begin{enumerate}[label=(\alph*)]
\item Every column of $P$ sums to one, while every column of $I$ sums
to one. Thus every column of $P-I$ sums to zero. Adding all other rows
to the bottom row produces that column-sum row, which is zero.
\item The sum of all original rows of $P-I$ is zero. The coefficients
of this relation are all $1$, not all zero, so the rows are dependent.
\item A dependent list of $n$ rows spans a space of dimension less than $n$.
Equivalently $\rank(P-I)<n$.
\item By rank--nullity, $\dim\Nul(P-I)=n-\rank(P-I)>0$.
Thus there is a nonzero $x$ satisfying $(P-I)x=0$, or $Px=x$.
\end{enumerate}
This proves existence of a nonzero fixed vector; by itself it does not
prove its entries can be chosen nonnegative and normalized, which is the
additional point mentioned in the exercise.

\exercise{4.9.18}{An explicit steady state for every two-state matrix}
For $P=\left(\begin{smallmatrix}1-\alpha&\beta\\\alpha&1-\beta\end{smallmatrix}\right)$,
the equation $Pq=q$ is $\alpha q_1=\beta q_2$.
If $\alpha+\beta>0$, choose
\[
 q=\left(\frac{\beta}{\alpha+\beta},
          \frac{\alpha}{\alpha+\beta}\right).
\]
Both entries are nonnegative and sum to one, and the balance equation
holds, so this is a steady-state vector. It is the unique normalized
solution because the balance and sum equations determine both entries.
If $\alpha=\beta=0$, then $P=I_2$. Every vector $(t,1-t)$ with
$0\le t\le1$ is a steady state; in particular $(1,0),(0,1)$ are two
independent ones. Thus the parenthetical claim means two independent
steady-state directions, not only two probability vectors in this case.

\exercise{4.9.19}{A row of ones records total probability}
\step{(a) Interpret the scalar $Sx$.}
For $S=[\,1\ \cdots\ 1\,]$, $Sx=\sum_i x_i$.
Therefore $x$ is a probability vector exactly when its entries are
nonnegative and $Sx=1$.
\step{(b) Interpret the entries of $SP$.}
Entry $j$ is the sum of column $j$ of $P$, which equals $1$.
Hence $SP=S$.
\step{(c) Preserve nonnegativity and the total separately.}
Every entry of $Px$ is a sum of products of nonnegative entries, so it
is nonnegative. Also $S(Px)=(SP)x=Sx=1$.
By part (a), $Px$ is a probability vector.

\exercise{4.9.20}{The square is stochastic}
Every column $p_j$ of $P$ is a probability vector. Column $j$ of $P^2$
is $Pp_j$, which is a probability vector by Exercise 4.9.19(c).
Thus each column of $P^2$ has nonnegative entries and sums to one,
which is exactly the definition of a stochastic matrix.

\exercise{4.9.21}{Observe convergence of all columns of powers}
\step{(a) Four powers of the supplied positive $4\times4$ matrix.}
Keep full internal precision in the products and round only the displayed
entries to four decimal places. The results are
\[ P^{2}=\begin{pmatrix}0.2779&0.2780&0.2803&0.2941\\0.3368&0.3355&0.3357&0.3335\\0.1847&0.1861&0.1833&0.1697\\0.2005&0.2004&0.2007&0.2027\end{pmatrix}.\]
\[ P^{3}=\begin{pmatrix}0.2817&0.2817&0.2817&0.2814\\0.3356&0.3356&0.3355&0.3352\\0.1817&0.1817&0.1819&0.1825\\0.2010&0.2010&0.2010&0.2009\end{pmatrix}.\]
\[ P^{4}=\begin{pmatrix}0.2816&0.2816&0.2816&0.2816\\0.3355&0.3355&0.3355&0.3355\\0.1819&0.1819&0.1819&0.1819\\0.2009&0.2009&0.2009&0.2009\end{pmatrix}.\]
\[ P^{5}=\begin{pmatrix}0.2816&0.2816&0.2816&0.2816\\0.3355&0.3355&0.3355&0.3355\\0.1819&0.1819&0.1819&0.1819\\0.2009&0.2009&0.2009&0.2009\end{pmatrix}.\]
Solving $(P-I)q=0$ together with $q_1+q_2+q_3+q_4=1$ gives
\[ q\approx\begin{pmatrix}0.2816326912\\0.3355052569\\0.181921411\\0.2009406409\end{pmatrix}.\]
The columns are rapidly becoming nearly equal to this fixed probability vector.
\step{(b) Powers of the more slowly mixing $3\times3$ matrix.}
\[ Q^{10}=\begin{pmatrix}0.8222&0.4044&0.5385\\0.0324&0.3966&0.1666\\0.1453&0.1990&0.2949\end{pmatrix}.\]
\[ Q^{20}=\begin{pmatrix}0.7674&0.6000&0.6690\\0.0637&0.2036&0.1326\\0.1688&0.1964&0.1984\end{pmatrix}.\]
\[ Q^{30}=\begin{pmatrix}0.7477&0.6815&0.7105\\0.0783&0.1329&0.1074\\0.1740&0.1856&0.1821\end{pmatrix}.\]
\[ Q^{40}=\begin{pmatrix}0.7401&0.7140&0.7257\\0.0843&0.1057&0.0960\\0.1756&0.1802&0.1783\end{pmatrix}.\]
\[ Q^{50}=\begin{pmatrix}0.7372&0.7269&0.7315\\0.0867&0.0951&0.0913\\0.1761&0.1780&0.1772\end{pmatrix}.\]
\[ Q^{60}=\begin{pmatrix}0.7360&0.7320&0.7338\\0.0876&0.0909&0.0894\\0.1763&0.1771&0.1767\end{pmatrix}.\]
\[ Q^{70}=\begin{pmatrix}0.7356&0.7340&0.7347\\0.0880&0.0893&0.0887\\0.1764&0.1767&0.1766\end{pmatrix}.\]
\[ Q^{80}=\begin{pmatrix}0.7354&0.7348&0.7351\\0.0881&0.0887&0.0884\\0.1764&0.1766&0.1765\end{pmatrix}.\]
Its second equilibrium equation gives $q_2=q_3/2$, and its first gives
$.03q_1=.05q_2+.10q_3$. Normalization then yields
\[
 q=(25/34,3/34,3/17)\approx(.73529412,.08823529,.17647059).
\]
The columns approach this vector more slowly, so an arbitrary fixed power
such as $100$ is not a universal certificate of four-place accuracy.
\step{(c) Explain the observation with Theorem 18.}
$P$ is positive, hence regular. $Q^2$ has all entries positive (the only
zero in $Q$ acquires a positive two-step route), so $Q$ is regular too.
Column $j$ of $P^k$ is $P^ke_j$, the evolving probability vector from the
initial state $e_j$. Theorem 18 says each such vector tends to the same
unique steady state. The same reasoning applies to $Q^k$ and explains the
observed column limits for any regular stochastic matrix.

\exercise{4.9.22}{A reproducible comparison of two numerical methods}
\goal Run both methods, report measured times, and distinguish numerical
approximation from a mathematical convergence guarantee.
A matrix program such as NumPy has no fixed largest permitted order;
the practical limit depends on memory and desired runtime. The following
experiment uses positive random matrices through order $1024$, the largest
order tested here, rather than claiming an intrinsic software size limit.
Normalize each random positive column to obtain a regular stochastic matrix.
For method 1 solve $(P-I)q=0$ with its final equation replaced by
$\sum_iq_i=1$. For method 2 use column one of $P^{100}$.

The following table records actual execution times in seconds on this
workspace, with one BLAS thread; input generation is outside the timing.
Each method is invoked by one matrix-program call after the shared setup.
Human typing time is user dependent and is not fabricated as an execution measurement.
\begin{center}\begin{tabular}{rrrr}\toprule
Order&Direct solve&Power $100$&Maximum difference\\\midrule
64&0.000149&0.000120&$6.09e-16$\\
128&0.000173&0.000829&$3.93e-16$\\
256&0.000788&0.005619&$5.89e-16$\\
512&0.004078&0.043638&$7.01e-16$\\
1024&0.028829&0.329893&$7.84e-16$\\
\bottomrule\end{tabular}\end{center}
Here is the complete experimental core; start Python with
\texttt{OPENBLAS\_NUM\_THREADS=1} for the stated thread setting.
\begin{verbatim}
import numpy as np
from time import perf_counter
rng = np.random.default_rng(4922)
for n in [64, 128, 256, 512, 1024]:
    P = rng.uniform(.1, 1, (n, n)); P /= P.sum(axis=0)
    A = P - np.eye(n); A[-1, :] = 1
    b = np.zeros(n); b[-1] = 1
    start = perf_counter()
    q = np.linalg.solve(A, b)
    direct = perf_counter() - start
    start = perf_counter()
    v = np.linalg.matrix_power(P, 100)[:, 0]
    power = perf_counter() - start
    print(n, direct, power, np.max(abs(q-v)))
\end{verbatim}
For these well-mixed positive samples, the two vectors agree to less than
$8\times10^{-16}$ in maximum coordinate error. This empirical accuracy
is not guaranteed for every regular stochastic matrix: a slowly mixing
matrix can require far more than $100$ steps, as Exercise 4.9.21 illustrates.
Solving the normalized linear system avoids choosing a mixing time and was
faster for the larger tested orders. Computing a power is easy to request
and can reveal transients, but a large power performs more matrix products
and still needs a convergence check. For finding only the steady state,
the direct normalized solve is preferable in this experiment.
